Input to the General-Head Boundary (GHB) Package is read from the file that has type ``GHB6'' in the Name File.  Any number of GHB Packages can be specified for a single groundwater flow model.

General-Head Boundary input can be specified using lists or grid arrays.  Grid array-based input for the GHB package is activated by providing READARRAYGRID within the OPTIONS block.  Instructions for specifying list-based general-head boundaries are provided in the previous section.

When grid array-based input is used for general-head boundaries, the DIMENSIONS block is optional.  The input array size is determined from the model grid.  Cells that define no data should contain the DNODATA (3.0E+30) value.  When defined, the MAXBOUND dimension represents the maximum number of defined general-head boundary (non-DNODATA) cells found in any stress period.  Defining MAXBOUND can reduce the memory footprint of the grid array-based GHB package, especially when the input grid arrays are sparse and the model grid is large.

\vspace{5mm}
\subsubsection{Structure of Blocks}
\vspace{5mm}

\noindent \textit{FOR EACH SIMULATION}
\lstinputlisting[style=blockdefinition]{./mf6ivar/tex/gwf-ghbg-options.dat}
\lstinputlisting[style=blockdefinition]{./mf6ivar/tex/gwf-ghbg-dimensions.dat}
\vspace{5mm}
\noindent \textit{FOR ANY STRESS PERIOD}
\lstinputlisting[style=blockdefinition]{./mf6ivar/tex/gwf-ghbg-period.dat}
\packageperioddescriptiongridarray{general-head boundary}

\vspace{5mm}
\subsubsection{Explanation of Variables}
\begin{description}
\input{./mf6ivar/tex/gwf-ghbg-desc.tex}
\end{description}

\vspace{5mm}
\subsubsection{Example Input File}
\lstinputlisting[style=inputfile]{./mf6ivar/examples/gwf-ghbg-example.dat}

\vspace{5mm}
\subsubsection{Available observation types}
General-Head Boundary Package observations include the simulated general-head boundary flow rates (\texttt{ghb}) and the general-head boundary discharge that is available for the MVR package (\texttt{to-mvr}). The data required for each GHB Package observation type is defined in table~\ref{table:gwf-ghbgobstype}. The sum of \texttt{ghb} and \texttt{to-mvr} is equal to the simulated general-head boundary flow rate. Negative and positive values for an observation represent a loss from and gain to the GWF model, respectively.

\begin{longtable}{p{2cm} p{2.75cm} p{2cm} p{1.25cm} p{7cm}}
\caption{Available GHB Package observation types} \tabularnewline

\hline
\hline
\textbf{Stress Package} & \textbf{Observation type} & \textbf{ID} & \textbf{ID2} & \textbf{Description} \\
\hline
\endhead

\hline
\endfoot

\input{../Common/gwf-ghbobs.tex}
\label{table:gwf-ghbgobstype}
\end{longtable}

\vspace{5mm}
\subsubsection{Example Observation Input File}
\lstinputlisting[style=inputfile]{./mf6ivar/examples/gwf-ghbg-example-obs.dat}
